\section{The sixteen-qubit exclusion}\label{sec:sixteen}

Assume that \(S\leq V_{16}\) defines a \(\qc{16,5,d\geq5}\) code.
Then \(|S|=2^{11}=2048\). The counts \(A_j\) include products of basis generators.
The argument has three steps. First we prove purity and restrict the
minimum stabilizer weight. Next we show that shortening on such a word
leaves one of four complete code classifications. Finally we test the
extension inequalities on every class in those lists.

\subsection{Purity and minimum stabilizer weight}

\begin{proposition}[Exact low-weight reduction]\label{prop:16-reduction}
Such a code is pure to distance five. Its minimum nonzero stabilizer
weight \(m\) belongs to \(\{6,7,8,9\}\).
\end{proposition}
\begin{proof}
The weight-enumerator equations count words of each weight without
specifying the full code. The MacWilliams identities express centralizer
counts in terms of stabilizer counts, and the shadow identities give
additional nonnegative counts. A certificate is an integer combination
of these necessary relations that gives a contradiction under a stated
case assumption. Its check uses integer arithmetic alone; the general
construction is given in Appendix~\ref{app:certificates}.

The six integer certificates in the repository exclude, in order:
\begin{enumerate}[label=(\alph*)]
\item a stabilizer word of weight one, two, or three;
\item a weight-four stabilizer word together with \(A_4\geq2\);
\item a weight-five stabilizer word under \(A_1=\cdots=A_4=0\);
\item \(A_1=\cdots=A_9=0\).
\end{enumerate}
For a proposed word of weight \(w\), local Clifford operations normalize
it to \(Z^w\). The split system counts complete cosets of
\(\langle Z^w\rangle\), with distance equations through weight four,
centralizer inclusion, and nonnegative shadow coefficients. Every actual
code yields a feasible point. Appendix~\ref{app:certificates} defines
the system and the exact arithmetic checks; the case assumptions and
negative certificate values are tabulated there.

It remains to exclude \(A_4=1\). Write \(\mathcal{K}_j(w)\) for the quaternary
Krawtchouk polynomial at length sixteen, and \(\operatorname{Sh}_j\)
for the shadow counts. Set
\[
 \Delta_0=\sum_wA_w=2048,\quad
 \Delta_j=\sum_w\mathcal{K}_j(w)A_w-2048A_j=0\ (1\leq j\leq4),
\]
and \(\Omega_j=\sum_w(-1)^w\mathcal{K}_j(w)A_w-2048\operatorname{Sh}_j=0\).
In the integer combination
\[
 -740\Delta_0-265\Delta_1-81\Delta_2-19\Delta_3-3\Delta_4
       +80\Omega_0+10\Omega_1-2\Omega_3,
\]
all coefficients vanish modulo \(4096\), except those of
\(A_1,A_2,A_3,A_4\), each congruent to \(2048\).
The right-hand side is \(-740\cdot2048\), which is divisible by
\(4096\).
\(A_1+A_2+A_3+A_4\) is even. The first certificates give
\(A_1=A_2=A_3=0\), \(A_4\leq1\), so \(A_4=0\).
Every centralizer word of weight below five must be a stabilizer;
therefore purity follows. The last two certificate cases give
\(6\leq m\leq9\).
\end{proof}

\subsection{The four shortened-code classifications}

Choose a weight-\(m\) word \(p\) with support \(T\), and normalize
\(p=Z^m\). We first need the \emph{entire} local stabilizer.

\begin{lemma}\label{lem:16-rankone}
For \(m\in\{6,7,8,9\}\), \(S\cap V_T=\langle p\rangle\).
\end{lemma}
\begin{proof}
If \(\dim(S\cap V_T)\geq2\), the shortened centralizer contains a
subcode with parameters
\[
 (10,2^{11},\geq5)_4,\quad(9,2^9,\geq5)_4,
       \quad\text{or}\quad(8,2^7,\geq5)_4
\]
for \(m=6,7,8\), respectively. No additive codes with these parameters exist
\cite{dodeca}; the repository also verifies their emptiness through the
complete lengthening census.
For \(m=9\), an independent local stabilizer word \(\zeta\) and \(\zeta+p\)
must both have weight nine by minimality of \(m\). \(\zeta\) has
\(X\) or \(Y\) at every site of \(T\). It anticommutes with
\(Z^9\) on an odd number of sites, contradicting isotropy.
\end{proof}

The shorter code \(C\) has length \(16-m\), dimension
\(22-2m\), and distance at least five. The quotient \(W/C\) has
dimension ten for every case. The complete classifications contain the following
numbers of physical equivalence classes:
\begin{center}
\begin{tabular}{@{}cccc@{}}\toprule
\(m\)&Shortened length&Binary dimension&Classes\\\midrule
6&10&10&37\\7&9&8&2,298\\8&8&6&1,579\\9&7&4&43\\\bottomrule
\end{tabular}
\end{center}
The four class counts are the exact entries of Table~1 in
Grassl et al. \cite{dodeca}.
The repository reconstructs representatives and checks their completeness
by a second physical graph encoding and orbit--stabilizer mass identities.
The classification procedure is specified in Section~\ref{sec:classification}.

\subsection{The finite prefix-map problem}

For \(R=\langle Z^m\rangle\), use the symplectic basis
\[
 Z_0,\ldots,Z_{m-2},\qquad X_iX_{m-1}\quad(0\leq i<m-1)
\]
of \(E\), with \(Z_{m-1}=\sum_{i<m-1}Z_i\) modulo \(R\).
Each basis image is a ten-bit syndrome. Every actual code gives a linear
map satisfying
\begin{equation}\label{eq:16-weight}
 \cost(e)+\delta(Me)\geq5\qquad(e\ne0).
\end{equation}
Only classes of minimum weight less than five constrain the search.
Omitting \(\Gamma=MJM^\tr\) enlarges the feasible set, so
nonexistence of a map in this larger set is sufficient.

The depth-first traversal assigns basis images and checks each newly
determined subset sum. The forward-checking version maintains every
remaining basis image's allowed domain. A value is removed only when a
specific already determined subset would violate \eqref{eq:16-weight}.
Every word is tested: each \(e\in E\) is a binary sum of the basis
vectors, and its image is the corresponding sum of assigned syndromes.
A Boolean satisfiability (SAT) encoding uses Boolean entries of \(M\), XOR chains
for its values, and a clause for every forbidden element/syndrome pair.
Each forbidden pair is a clause
requiring the entries of \(M\) to avoid that particular map value.
The two encodings construct their coset labels separately.

\begin{lemma}[Physical normalization for the weight-only search]\label{lem:16-symmetry}
If a map satisfying \eqref{eq:16-weight} exists, one exists with
\[
 \xi_0<\xi_1<\cdots<\xi_{m-1},\qquad
 \eta_i<\eta_i+\xi_i\quad(0\leq i<m-1),
\]
where \(\xi_i=M(Z_i)\), \(\eta_i=M(X_iX_{m-1})\), and labels use
any fixed total order.
\end{lemma}
\begin{proof}
The singleton and pair constraints force all \(\xi_i\) to be distinct
and nonzero. A physical permutation of the \(m\) sites fixes \(Z^m\)
and preserves coset minimum weights, so orders the singleton images.
After this permutation, the map is re-expressed in the displayed standard
basis; the dependent image is still
\(\xi_{m-1}=\sum_{i=0}^{m-2}\xi_i\).
The phase operation
\(X\leftrightarrow Y\) on site \(i<m-1\) fixes every \(Z_j\),
fixes the other complementary basis vectors, and replaces
\(X_iX_{m-1}\) by \(X_iX_{m-1}+Z_i\). It exchanges
\(\eta_i\) and \(\eta_i+\xi_i\) independently for each \(i\).
These physical operations preserve all weight inequalities.
\end{proof}

\begin{proposition}[Complete sixteen-qubit extension exclusion]\label{prop:16-search}
None of the 3,957 shortened-code classes admits a map satisfying
\eqref{eq:16-weight}.
\end{proposition}
\begin{proof}[Computer-assisted proof]
The recorded reports cover completed searches for every class.
For \(m\leq8\), a symmetry-free depth-first search and a separately
constructed SAT encoding both exclude every class. For \(m=9\), two
depth-first implementations use Lemma~\ref{lem:16-symmetry}, with
different coset coordinates and search orders. Both exclude all 43 classes.
A known pure \(\qc{16,4,5}\) code is accepted by
both encodings, with and without their symmetry restrictions.
Appendix~\ref{app:search} specifies the finite predicate and completeness
argument; Appendix~\ref{app:evidence} identifies the evidence and its checks.
\end{proof}

Proposition~\ref{prop:16-reduction}, Lemma~\ref{lem:16-rankone}, and
Theorem~\ref{thm:extension} put every hypothetical code into one of these
weight-admissible extension problems. Proposition~\ref{prop:16-search}
excludes them all. This proves the first assertion of
Theorem~\ref{thm:main}.
